\section{Length-two Witt arithmetic for divided Frobenius}
\label{sec:witt-two-arithmetic}

Let $p$ be prime, let $k$ be a perfect field of
characteristic $p$, and let $W_2(k)$ be the actual ring
of length-two $p$-typical Witt vectors. Coordinates
are denoted $(x_0,x_1)$. These elementary scalar facts
supply the arithmetic for divided Frobenius maps;
they do not assert existence of a scheme lift or a
Deligne--Illusie splitting.

\begin{definition}[The original reduction map]
\label{def:witt-two-reduction}
\leanok
\lean{EnvelopeProperty.DeligneIllusie.WittTwo.reduction,EnvelopeProperty.DeligneIllusie.WittTwo.reduction_apply,EnvelopeProperty.DeligneIllusie.WittTwo.reduction_surjective}
The first Witt coordinate defines a surjective ring
homomorphism $\rho:W_2(k)\to k$.
\end{definition}
\begin{proof}
The constant coefficient of a full Witt vector is a
ring homomorphism. It vanishes on the kernel of
truncation to two coordinates, because a vector in
that kernel has zero first coordinate. It therefore
factors through the surjective truncation map. The
result sends a two-coordinate vector to its first
coordinate. For $a\in k$, the truncated Teichmuller
vector $[a]$ has first coordinate $a$, proving
surjectivity.
\end{proof}

\begin{lemma}[Multiplication by the prime]
\label{lem:witt-two-prime-coordinates}
\leanok
\lean{EnvelopeProperty.DeligneIllusie.WittTwo.mul_prime_coeff_zero,EnvelopeProperty.DeligneIllusie.WittTwo.mul_prime_coeff_one}
For every $x\in W_2(k)$, the actual Witt product
$x p$ has coordinates $(0,x_0^p)$.
\uses{def:witt-two-reduction}
\end{lemma}
\begin{proof}
Choose a full Witt vector truncating to $x$. In
characteristic $p$, the coordinates of its product
with $p$ are zero in degree zero and the $p$th power
of the preceding coordinate in each positive degree.
Truncation is a ring homomorphism and preserves the
first two coordinates, giving the assertion.
\end{proof}

\begin{lemma}[Square-zero prime and its annihilator]
\label{lem:witt-two-prime-annihilator}
\leanok
\lean{EnvelopeProperty.DeligneIllusie.WittTwo.mul_prime_eq_zero_iff,EnvelopeProperty.DeligneIllusie.WittTwo.prime_square_zero}
For $x\in W_2(k)$, one has $xp=0$ if and only if
$\rho(x)=0$. In particular, $p^2=0$ in $W_2(k)$.
\uses{lem:witt-two-prime-coordinates}
\end{lemma}
\begin{proof}
The product formula gives $xp=0$ if and only if
$x_0^p=0$. As $k$ is a field and $p$ is positive,
this is equivalent to $x_0=0$, which is
$\rho(x)=0$. Since $\rho$ is a ring homomorphism
and $k$ has characteristic $p$, it sends $p$ to
zero. Applying the first assertion to $x=p$
gives $p^2=0$.
\end{proof}

\begin{lemma}[The reduction kernel is exactly the prime ideal]
\label{lem:witt-two-prime-kernel}
\leanok
\lean{EnvelopeProperty.DeligneIllusie.WittTwo.exists_mul_prime_iff}
An element $x$ of the same Witt ring is a multiple
of $p$ if and only if $\rho(x)=0$.
\uses{lem:witt-two-prime-coordinates,lem:witt-two-prime-annihilator}
\end{lemma}
\begin{proof}
If $x=yp$, its first coordinate is zero by the
product formula. Conversely suppose $x_0=0$.
Perfectness gives $b\in k$ with $b^p=x_1$.
Take $y$ to be the truncated Teichmuller vector
$[b]$. The product $yp$ has coordinates
$(0,b^p)=(x_0,x_1)$, so it equals $x$.
\end{proof}

\begin{lemma}[Uniqueness of division after reduction]
\label{lem:witt-two-divided-prime-uniqueness}
\leanok
\lean{EnvelopeProperty.DeligneIllusie.WittTwo.mul_prime_eq_iff}
For $x,y\in W_2(k)$, the equality $xp=yp$ holds
if and only if $\rho(x)=\rho(y)$. Thus dividing
an element of $pW_2(k)$ by $p$ determines a unique
residue in $k$.
\uses{lem:witt-two-prime-annihilator,lem:witt-two-prime-kernel}
\end{lemma}
\begin{proof}
Distributivity identifies $xp-yp$ with $(x-y)p$.
The annihilator formula says that this vanishes
if and only if $\rho(x-y)=0$. Additivity of
$\rho$ identifies the latter equality with
$\rho(x)=\rho(y)$. Existence of a division for
every element of $pW_2(k)$ is the preceding lemma.
\end{proof}
